
\textbf{Long-context LLMs.}
Several lines of work have improved the context length of LLMs. For instance, more efficient transformer architectures via sparsifying the attention \citep{child2019generating, beltagy2020longformer}, low-rank approximations \citep{wang2020linformer}, and neural memory \cite{lee2019set}. Another line of work aims to extend context windows beyond the length they were original trained for, their training size, such as \citet{alibi, chen2023extending}. \ours builds upon these improvements in context length as they improve the size of the main memory in \ours. Our main contribution is a hierarchical tiered memory that uses a long-context LLM as the implementation of main memory.

\textbf{Retrieval-Augmented Models.}
The design of the external memory of \ours builds upon much prior work augmenting LLMs with relevant inputs from external retrievers ~\cite{ram2023context, retro, fb-dense-retrieval, lewis2020retrieval, guu2020retrieval, radit}. In particular, \citet{jiang2023active} propose FLARE, a method that allows the LLM to actively decide when and what to retrieve during the course of generation.  \citet{Trivedi2022InterleavingRW} interleave retrieval with Chain-of-Thoughts reasoning to improve multi-step question answering. 


\textbf{LLMs as agents.}
Recent work has explored augmenting LLMs with additional capabilities to act as agents in interactive environments. 
\citet{park2023simulacra} propose adding memory to LLMs and using the LLM as a planner, and observe emergent social behaviors in a multi-agent sandbox environment (inspired by \emph{The Sims} video game) where agents can perform basic activities such as doing chores/hobbies, going to work, and conversing with other agents.
\citet{nakano2021webgpt} train models to search the web before answering questions, and use similar pagination concepts to \ours to control the underlying context size in their web-browsing environment. 
\citet{yao2022react} showed that interleaving chain-of-thought reasoning \citep{wei2022chainofthought} can further improve the planning ability of interactive LLM-based agents; similarly in \ours, LLM is able to `plan out loud' when executing functions. %
\citet{liu2023agentbench} introduced a suite of LLM-as-an-agent benchmarks to evaluate LLMs in interactive environments, including video games, thinking puzzles, and web shopping. 
In contrast, our work focuses on tackling the problem of equipping agents with long-term memory of user inputs.


